\documentclass[tikz,border=8pt]{standalone}
\usepackage{tikz}
\usepackage{amsmath}
\usepackage{pgfplots}
\pgfplotsset{compat=1.18}
\usepackage{ctex}
\begin{document}
\begin{tikzpicture}
\begin{axis}[
  width=10cm, height=7cm,
  xmin=-3.5, xmax=3.5,
  ymin=-0.05, ymax=1.1,
  xlabel={$x$},
  ylabel={$\Phi(x)$},
  title={\textbf{标准正态累积分布函数 $\Phi(x)$}},
  axis lines=left,
  xtick={-3,-2,-1,0,1,2,3},
  ytick={0,0.1,0.2,0.3,0.4,0.5,0.6,0.7,0.8,0.9,1.0},
  yticklabel style={font=\scriptsize},
  xticklabel style={font=\scriptsize},
  grid=major, grid style={gray!25},
  clip=false
]

  % CDF curve using tabulated values of standard normal
  \addplot[blue!70!black, ultra thick] coordinates {
    (-3.5, 0.0002) (-3.0, 0.0013) (-2.5, 0.0062) (-2.0, 0.0228)
    (-1.96, 0.0250) (-1.645, 0.0500) (-1.5, 0.0668) (-1.0, 0.1587)
    (-0.5, 0.3085) (0.0, 0.5000) (0.5, 0.6915) (1.0, 0.8413)
    (1.5, 0.9332) (1.645, 0.9500) (1.96, 0.9750) (2.0, 0.9772)
    (2.5, 0.9938) (3.0, 0.9987) (3.5, 0.9998)
  };

  % Key points with dashed lines
  % Phi(0) = 0.5
  \addplot[dashed, gray] coordinates {(0,0) (0,0.5)};
  \addplot[dashed, gray] coordinates {(-3.5,0.5) (0,0.5)};
  \fill[red] (axis cs:0,0.5) circle (2.5pt);
  \node[font=\scriptsize, red, right] at (axis cs:0.05,0.52) {$\Phi(0)=0.5$};

  % Phi(1.96) = 0.975
  \addplot[dashed, orange!80] coordinates {(1.96,0) (1.96,0.975)};
  \addplot[dashed, orange!80] coordinates {(-3.5,0.975) (1.96,0.975)};
  \fill[orange] (axis cs:1.96,0.975) circle (2.5pt);
  \node[font=\scriptsize, orange!80!black, left] at (axis cs:-0.1,0.975)
    {$0.975$};
  \node[font=\scriptsize, below] at (axis cs:1.96,-0.02) {$1.96$};

  % Symmetry annotation
  \draw[->, green!60!black, thin] (axis cs:-2.0,0.35) -- (axis cs:-1.1,0.16);
  \node[font=\scriptsize, green!50!black, align=center] at (axis cs:-2.3,0.38)
    {$\Phi(-x)=1-\Phi(x)$};

  % Horizontal asymptotes
  \addplot[dotted, gray, thin] coordinates {(-3.5,1.0) (3.5,1.0)};
  \node[font=\scriptsize, gray, right] at (axis cs:3.3,1.02) {$1$};

\end{axis}
\end{tikzpicture}
\end{document}
